% Pro P. Sulla (pro-sulla) --- English translation
% Translated by Claude (Anthropic) from the Latin (Perseus).
% Style guide: STYLE.md.

\ciceroSection{1} What I should most have wished, judges, was that Publius Sulla had been able to keep the splendour of his standing as it once was, and after the disaster he suffered to gather some fruit from his moderation. But since hostile chance has fallen out so that at the height of his honour he was overthrown by the common odium that attaches to canvassing and by the singular hatred of Autronius, and since amid these wretched and battered remains of his former fortune he still has men whose temper not even his ruin could glut --- though I take great heaviness of mind from his troubles, yet among my other griefs I bear with patience that an occasion has been given me on which good men may recognize a mildness and mercy in me once familiar to them all, now as it were laid aside, and on which the wicked and ruined citizens, mastered and beaten, will admit that with the commonwealth tottering I was vehement and brave, and with the commonwealth saved gentle and merciful.

\ciceroSection{2} And since Lucius Torquatus, my close friend and connection, judges, has thought that, if he should violate that intimacy and friendship of ours in his prosecution, he could detract something from the authority of my defence, I shall couple with the warding off of this man's danger the defence of my own duty. I should certainly not be using this kind of speech now, judges, were the matter mine alone --- for in many places the chance has been given me, and will often be given again, of speaking to my own praise. But just as he saw that as much as he stripped from my authority, by so much would he diminish this man's protection, so I, on my side, judge that if I prove to you the reasoning of my act and the consistency of this duty and defence, I shall also prove the case of Publius Sulla.

\ciceroSection{3} And first, Lucius Torquatus, this is what I ask of you: why do you set me apart from the rest of the most distinguished men, the leading citizens, in this duty and in the right of defence? What is the reason why what Quintus Hortensius has done, that man most distinguished and most adorned, is not censured by you, but mine is? For if a plot was formed by Publius Sulla to set this city on fire, to extinguish our empire, to wipe out our state, these things should bring me greater grief than they should bring Hortensius, greater hatred toward me; my judgment finally as to who is to be helped in such cases, who attacked, who defended, who deserted, ought to be the weightier. ``Just so,'' he says; ``for it was you who tracked down, you who exposed the conspiracy.''

\ciceroSection{4} When he says this, he does not consider that the man who exposed it took care that all should see what before had been hidden. So that conspiracy, if it has been exposed through me, is as exposed to Hortensius as it is to me. And when you see him --- a man endowed with such honour, authority, virtue, and judgment --- not hesitating to defend the innocence of Publius Sulla, I ask why the path to the case that lay open to Hortensius should have been shut off to me. I ask this too: if you think that I, who am defending him, should be censured, what then do you think of these greatest of men, these most distinguished citizens, by whose zeal and dignity you see this trial honoured, this case adorned, this man's innocence defended? For there is not just one mode of defence, the one set out in a speech: all who attend, who are anxious, who wish him saved, defend him each in his own measure of standing.

\ciceroSection{5} Or was I, when I saw on these benches the ornaments and lights of the commonwealth, to be unwilling to appear among them --- I, who climbed up to that very place, to this most exalted seat of dignity and honour, by my own great labours and dangers? And so that you may understand, Torquatus, whom you accuse, in case it offends you that I, who have defended no one in this kind of trial, do not fail Publius Sulla, recall who else are present in his support: you will see that my judgment in his case and in the others, and theirs, has been one and the same.

\ciceroSection{6} Which of us stood by Vargunteius? No one --- not even this Quintus Hortensius, who alone had previously defended him on a charge of canvassing. For he no longer thought himself bound by any duty to that man, since by committing so great a crime he had broken the bond of all duties. Which of us thought that Servius Sulla, that Publius, that Marcus Laeca, that Gaius Cornelius should be defended? Which of us stood by them? No one. Why? Because in other cases good men do not think that even guilty men, if they are connections, should be deserted; in this charge, however, the fault is not only one of fickleness but there is even a sort of contagion of crime, if you defend a man whom you suspect of being bound by parricide of the fatherland.

\ciceroSection{7} What of Autronius? Did not his comrades, his colleagues, his old friends, of whom he once had had abundance, did not all the leading men of the commonwealth fail him? Indeed many of them even injured him by their evidence. They had judged that the wickedness was such as not only ought not to be hidden by them, but rather to be opened up and brought into the light. So why is there anything to wonder at if you see me present in this case among the same men with whom you understand I was absent in the others? Unless, indeed, you wish me alone, beyond the rest, to be thought savage, harsh, inhuman, endowed with a singular barbarity and cruelty.

\ciceroSection{8} If you fix this character on me for the whole of my life, Torquatus, on account of my actions, you are very wrong indeed. By nature I am merciful, by my country's call severe; cruel, neither my country nor my nature would have me. The very vehement and harsh character which the time and the commonwealth then forced on me, my own will and very nature have already laid aside. For that called for severity for a brief time; this calls for mercy and mildness for the whole of life.

\ciceroSection{9} There is therefore no reason for you to drag me alone away out of so great a company of the most distinguished men. The duty is single, and the cause of all good men one and the same. There will be nothing for you to wonder at hereafter if in that quarter where you have noticed these you will see me also. For there is no cause of mine private to me in the commonwealth: a time for action was more peculiarly mine than the rest's, but the cause of grief and fear and danger was for all of us in common. Nor could I then have been the leader of our salvation, had others not been willing to be my companions. So what was peculiarly mine as consul, beyond the others, must now, as a private man, be common to me with the rest. Nor do I say this to share the odium, but to share the praise: I impart to no one a part of my burden, to all good men a share of the glory.

\ciceroSection{10} ``You gave evidence against Autronius,'' he says; ``you defend Sulla.'' This whole charge is of the kind, judges, that, if I am inconstant and frivolous, neither should credit have been given to my evidence nor authority to my defence; but if there is in me regard for the commonwealth, scruple in private duty, zeal for keeping the goodwill of good men, the prosecutor's last word should be that Sulla is being defended by me, that Autronius was injured by my testimony. For I see now that I bring not only zeal for defending cases, but some weight of opinion and authority too. This I shall use moderately, judges, and should not be using at all if he had not driven me to it.

\ciceroSection{11} Two conspiracies, Torquatus, are set up by you: one which is said to have been formed in the consulship of Lepidus and Volcatius, when your father was consul-designate; the other in mine. In each of these you say Sulla was involved. As to your father's, that bravest of men and best of consuls, you know I had no part in his counsels; you know that, although I was on the most familiar terms with you, yet I had no share in those proceedings or those discussions --- because, I take it, I was not yet deeply involved in public life, because I had not yet attained the goal of office I had set myself, because canvassing and forensic work drew me away from all such reflection.

\ciceroSection{12} Who, then, took part in your counsels? All these men whom you see attending here, and Quintus Hortensius first --- he, both because of his standing and dignity and uncommon spirit toward the commonwealth, and because of his close friendship and high love for your father, was moved as much by the dangers common to all as by the singular dangers of your father. So the charge of that conspiracy has been answered by the man who was present, who knew, who shared your counsel and your fear; and although his speech in beating off this charge was a most copious and most ornate one, yet there was no less authority in it than there was skill. As for that conspiracy, then, which is said to have been formed against you, reported to you, brought to light by you, I could not be a witness: not only have I learned nothing of it firsthand, but the rumour of that suspicion scarcely reached my ears.

\ciceroSection{13} The men who shared your counsel, who learned of these things together with you, who themselves were thought to be the targets at the time, who did not stand by Autronius, who gave grave evidence against him --- these defend this man, stand by him, declare in his danger that they were deterred not by a charge of conspiracy (so that they should be absent from the others) but by the wickedness of those men. The time and charge of my consulship, however, of the greatest conspiracy, will be defended by me. And this division of the defence between us has not been made fortuitously, judges, nor at random; but when we saw that we were being called as patrons in those charges of which we could be witnesses, each of us thought that he should take on himself the part on which he himself could know and judge something.

\ciceroSection{14} And since you have heard Hortensius diligently on the charges of the earlier conspiracy, attend first to this on the conspiracy that was formed in my consulship. Many things, when I was consul, did I hear about the supreme dangers of the commonwealth, many things did I investigate, many things did I learn; never any messenger about Sulla reached me, no informer, no letter, no suspicion. Much weight perhaps that voice ought to carry --- the voice of the man who as consul tracked down by his own counsel the plots against the commonwealth, exposed them by truth, avenged them by greatness of spirit --- when he says that he heard nothing about Publius Sulla, suspected nothing. But I do not yet use this voice for defending him: I shall use it rather to clear myself, that Torquatus may stop wondering that I, who did not stand by Autronius, defend Sulla.

\ciceroSection{15} For what was Autronius's case, what is Sulla's? He sought to break up the canvassing trial first by stirring up a riot of gladiators and runaway slaves, then --- as we all saw --- by stoning and a rush of mob. Sulla, when his own modesty and dignity were of no help to him, sought no resort. The other, once condemned, conducted himself --- not only in his counsels and discourse, but in his very face and look --- as a man manifestly hostile to the highest orders, dangerous to all good men, an enemy to his country. This man thought himself so broken and beaten by that disaster that he reckoned nothing remained to him of his former dignity except what he had kept by his own modesty.

\ciceroSection{16} In this present conspiracy, what was so close as that man's bond with Catiline, with Lentulus? What partnership in the best things has any men joined together so close as theirs in crime, in lust, in audacity? What infamy did Lentulus conceive without Autronius? What deed did Catiline commit without that same man? Meanwhile Sulla with these men sought neither night nor solitude, but did not even join them in ordinary speech and meeting.

\ciceroSection{17} Him the Allobroges, those most truthful informers in the gravest of matters, him the letters and messages of many men convicted; Sulla meanwhile no one accused, no one named. Lastly, when Catiline had been driven out --- or had been let out --- of the city, that man sent arms, horns, trumpets, fasces, standards, legions; left behind within the walls, awaited from outside, hemmed in only by the punishment of Lentulus, he turned at last to fear, never to sanity. This man on the contrary kept so quiet that for the whole of that time he was at Naples --- where men involved in this kind of suspicion are not thought to have been, and where the place itself is more suited to consoling broken spirits than to inflaming them. Therefore, on account of this great unlikeness of men and of cases, I have shown myself unlike in each.

\ciceroSection{18} For Autronius used to come to me, and used to come often, a suppliant with many tears begging that I would defend him, and recalled that he had been my schoolfellow in boyhood, my familiar in youth, my colleague in the quaestorship. He brought up many services of mine to him, some too of his own to me. By all this, judges, I was so bent and broken in spirit that already I was laying aside the memory of the plots he had made against me, was already forgetting that Gaius Cornelius had been sent by him to butcher me in my own house, in sight of my wife and children. Had he plotted these things against me alone, soft and gentle as I am of disposition, by Hercules I should never have held out against his tears and prayers; but when into my mind there came

\ciceroSection{19} the danger of my country, of you all, of this city, of those shrines and temples, of children at the breast, of matrons and maidens, and when those hostile and deadly torches and the universal fire of the whole city, when the weapons, when the slaughter, when the blood of citizens, when the ash of my country began to play before my eyes and to chafe my mind by memory, then at last I withstood him --- and not only that enemy and parricide but also those connections of his, the Marcelli, father and son, of whom the one held in my eyes the gravity of a parent, the other the sweetness of a son. Nor did I think it possible without supreme wickedness to defend in another's accomplice, when I knew it, the same crime which I had punished in others.

\ciceroSection{20} And I, the same man, could not endure Publius Sulla as a suppliant, nor look upon those same Marcelli weeping for this man's dangers, nor sustain the prayers of his connection Marcus Messala. For neither was the case opposed to nature, nor did either man or matter fight against my mercy. Nowhere had there been a name, nowhere a trace, no charge, no informer's word, no suspicion. I took the case up, Torquatus, I took it up, and gladly --- so that I, whom good men, as I hope, have always thought consistent, should not even by the wicked be called cruel.
